\chapter{Συμπεράσματα και Μελλοντικές Επεκτάσεις}
\label{ch:conclusions}

\section{Σύνοψη ευρημάτων}
\label{sec:conclusions-findings}

Τα αποτελέσματα της εργασίας δείχνουν αρχικά ότι η απόσταξη γνώσης
αποτελεί σαφώς αποτελεσματικότερη προσέγγιση από την επιβλεπόμενη
εκπαίδευση των μοντέλων-\tl{student} από την αρχή. Όταν οι
αρχιτεκτονικές \tl{V-JEPA Base}, \tl{Small} και \tl{Tiny} εκπαιδεύονται
αποκλειστικά με τις ετικέτες του \tl{HMDB51}, η ακρίβειά τους παραμένει
σε πολύ χαμηλά επίπεδα, καθώς το περιορισμένο μέγεθος του συνόλου δεν
επαρκεί για την αποτελεσματική εκπαίδευση των αντίστοιχων
\tl{Vision Transformer backbones}. Αντίθετα, η καθοδήγησή τους από το
προεκπαιδευμένο μοντέλο-\tl{teacher} επιτρέπει την αξιοποίηση της γνώσης
που αυτό απέκτησε κατά την αυτοεπιβλεπόμενη προεκπαίδευσή του και
οδηγεί σε σημαντικά υψηλότερες επιδόσεις.

Καθοριστικός παράγοντας για την αποτελεσματικότητα της απόσταξης
αποδείχθηκε το είδος της πληροφορίας που χρησιμοποιείται ως στόχος
εκπαίδευσης. Η απόσταξη μόνο των τελικών προβλέψεων προσέφερε περιορισμένη
βελτίωση, ενώ η προσθήκη των αναπαραστάσεων του \tl{backbone} στην
υβριδική στρατηγική αύξησε την ακρίβεια, χωρίς όμως να την
προσεγγίσει σε εκείνη του μοντέλου-\tl{teacher}. Η ουσιαστική
βελτίωση προέκυψε μόνο όταν η απόσταξη περιορίστηκε αποκλειστικά στις
αναπαραστάσεις, με την εκπαίδευση του \tl{attentive classifier} να
πραγματοποιείται σε ξεχωριστό στάδιο. Το αποτέλεσμα αυτό δείχνει ότι, για ένα
μοντέλο εκμάθησης αναπαραστάσεων όπως το \tl{V-JEPA}, οι έξοδοι του
\tl{backbone} παρέχουν ουσιαστικότερη καθοδήγηση από τις τελικές
προβλέψεις, επιτρέποντας στον \tl{student} να μάθει χρήσιμα
οπτικοχρονικά χαρακτηριστικά.

Ως προς τα δεδομένα, η απόσταξη στο \tl{UCF101} οδήγησε σε επίδοση
παρόμοια με εκείνη στο \tl{HMDB51}, παρότι η εκπαίδευση του
\tl{attentive classifier} και η τελική αξιολόγηση πραγματοποιήθηκαν
στο δεύτερο. Τα δύο σύνολα δεν ακολουθούν την ίδια κατανομή, καθώς
διαφέρουν ως προς την προέλευση των βίντεο και έχουν περιορισμένη
επικάλυψη κλάσεων (Ενότητα~\ref{sec:dataset-comparison}), οπότε το
\tl{UCF101} εισάγει μια μετατόπιση δεδομένων (\tl{dataset drift}) σε
σχέση με την εργασία-στόχο. Παρά τη μετατόπιση αυτή, τα αποτελέσματα
έδειξαν ότι, στο πλαίσιο των συνόλων που εξετάστηκαν, τα δεδομένα
απόσταξης δεν είναι απαραίτητο να προέρχονται από το σύνολο της
εργασίας-στόχου, εφόσον έχουν παρόμοιο σημασιολογικό περιεχόμενο. Καθώς η απόσταξη του \tl{backbone} δεν
απαιτεί ετικέτες, μπορεί επιπλέον να εφαρμοστεί σε μη επισημασμένα
σύνολα βίντεο σχετικής θεματολογίας. 

Τόσο η αύξηση του αριθμού των εποχών απόσταξης όσο και η χρήση
μεγαλύτερου όγκου δεδομένων συμβάλλουν στη βελτίωση της επίδοσης. Οι υψηλότερες επιδόσεις προέκυψαν από τον συνδυασμό των
δύο συνόλων, με το όφελος να είναι εντονότερο για τις μικρότερες
αρχιτεκτονικές, γεγονός που υποδεικνύει τη σημασία του πλήθους και της
ποικιλίας των δεδομένων.

Με την καλύτερη πειραματική διαμόρφωση, το \tl{V-JEPA Base} έφτασε σε
ακρίβεια \(76{,}60\%\), έναντι \(81{,}31\%\) του
\tl{V-JEPA Large}, διατηρώντας το \(94{,}2\%\) της επίδοσής του με
\(70{,}3\%\) λιγότερες παραμέτρους. Το \tl{V-JEPA Small} διατήρησε το
\(89{,}5\%\) της επίδοσης με \(92{,}4\%\) λιγότερες παραμέτρους, ενώ το
\tl{V-JEPA Tiny} διατήρησε το \(75{,}4\%\) με μείωση παραμέτρων κατά
\(98\%\). Παράλληλα, μειώθηκαν σημαντικά οι απαιτήσεις μνήμης και
επιτεύχθηκε έως και \(4{,}2\) φορές μεγαλύτερο \tl{throughput}.

Συνολικά, τα αποτελέσματα απαντούν θετικά στο βασικό ερευνητικό ερώτημα
της εργασίας: η γνώση ενός μεγάλου αυτοεπιβλεπόμενου μοντέλου βίντεο
μπορεί να μεταφερθεί αποτελεσματικά σε σημαντικά μικρότερες
αρχιτεκτονικές, διατηρώντας μεγάλο μέρος της επίδοσής του και μειώνοντας
παράλληλα τις υπολογιστικές απαιτήσεις. Η επιλογή του κατάλληλου
μοντέλου εξαρτάται από τις ανάγκες της εφαρμογής. Το
\tl{V-JEPA Base} προσφέρει επίδοση κοντά σε εκείνη του
μοντέλου-\tl{teacher}, το \tl{V-JEPA Small} προσφέρει έναν καλό συμβιβασμό μεταξύ επίδοσης και υπολογιστικού κόστους, ενώ το
\tl{V-JEPA Tiny} μπορεί να αξιοποιηθεί σε περιβάλλοντα με αυστηρότερους
περιορισμούς πόρων, όπου είναι αποδεκτή μεγαλύτερη μείωση της
ακρίβειας. Επομένως, η απόσταξη αναπαραστάσεων αποτελεί μια υποσχόμενη προσέγγιση
για τη μεταφορά της γνώσης μεγάλων μοντέλων θεμελίωσης βίντεο σε
μικρότερες και υπολογιστικά οικονομικότερες αρχιτεκτονικές,
διευρύνοντας τη δυνατότητα αξιοποίησής τους σε εφαρμογές και
περιβάλλοντα όπου οι διαθέσιμοι πόροι δεν επαρκούν για την εκτέλεση των
αρχικών μοντέλων.

\section{Περιορισμοί}
\label{sec:conclusions-limitations}

Τα συμπεράσματα της εργασίας πρέπει να ερμηνεύονται λαμβάνοντας υπόψη
ορισμένους περιορισμούς. Αρχικά, η αξιολόγηση πραγματοποιήθηκε σε μία μόνο
εργασία, την αναγνώριση ανθρώπινων ενεργειών, και στον πρώτο επίσημο
διαχωρισμό του \tl{HMDB51}. Επιπλέον, κάθε πείραμα εκτελέστηκε μία φορά με το ίδιο \tl{random seed},
χωρίς επανάληψη με διαφορετικές αρχικοποιήσεις. Τεχνικές επικύρωσης όπως
η \tl{k-fold cross-validation} και η \tl{leave-one-subject-out (LOSO)}
δεν εφαρμόστηκαν, καθώς απαιτούν την επανάληψη ολόκληρης της εκπαίδευσης
για κάθε διαχωρισμό, αυξάνοντας σημαντικά τον ήδη μεγάλο χρόνο της
απόσταξης γνώσης. Επομένως, μικρές
διαφορές μεταξύ των πειραματικών διαμορφώσεων, της τάξης της μίας
ποσοστιαίας μονάδας, δεν μπορούν να θεωρηθούν
απολύτως αξιόπιστες.

Λόγω των διαθέσιμων υπολογιστικών πόρων, η απόσταξη πραγματοποιήθηκε σε
σύνολα δεδομένων σχετικά μικρής κλίμακας και η διερεύνηση των
υπερπαραμέτρων δεν ήταν εξαντλητική. Ορισμένες διαμορφώσεις, όπως οι
διαφορετικοί προβολείς της υβριδικής στρατηγικής, εξετάστηκαν μόνο στο
\tl{V-JEPA Base} και δεν εφαρμόστηκαν σε όλες τις αρχιτεκτονικές.

Ακόμη, η σύγκριση των μοντέλων περιορίστηκε στις τρεις στρατηγικές
απόσταξης, το μοντέλο-\tl{teacher} και τις αντίστοιχες αρχιτεκτονικές
εκπαιδευμένες από την αρχή, χωρίς άμεση σύγκριση με άλλες τεχνικές
συμπίεσης. Παράλληλα, η μελέτη περιορίστηκε στην οικογένεια
\tl{V-JEPA}, οπότε δεν εξετάστηκε κατά πόσο τα συμπεράσματα
γενικεύονται σε διαφορετικές οικογένειες μοντέλων βίντεο.

Το κόστος της ίδιας της απόσταξης παραμένει υψηλό, καθώς σε κάθε βήμα
εκτελείται και το σημαντικά μεγαλύτερο μοντέλο-\tl{teacher}. Για τον
λόγο αυτό, ο χρόνος εκπαίδευσης δεν μειώνεται σημαντικά για τα μικρότερα
μοντέλα-\tl{student}. Ενδεικτικά, η απόσταξη για \(200\) εποχές στον
πλήρη συνδυασμό των \tl{HMDB51} και \tl{UCF101} απαίτησε περίπου
\(18\) για καθένα από τα τρία μοντέλα.
Το κόστος αυτό εξηγεί σε μεγάλο βαθμό και τους περιορισμούς που
αναφέρθηκαν παραπάνω ως προς τον αριθμό των επαναλήψεων και την
έκταση της διερεύνησης.

Τέλος, οι μετρήσεις της μέγιστης \tl{VRAM} και του \tl{throughput}
πραγματοποιήθηκαν αποκλειστικά σε \tl{GPU NVIDIA A100}, με συγκεκριμένο
μέγεθος παρτίδας (\tl{batch size}) και αριθμητική ακρίβεια. Συνεπώς, τα σχετικά
αποτελέσματα εκφράζουν τη αποδοτικότητα των μοντέλων στις
συγκεκριμένες συνθήκες και δεν μεταφέρονται απαραίτητα αυτούσια σε
διαφορετικό \tl{hardware}. Η καθυστέρηση (\tl{latency}) ανά βίντεο, η κατανάλωση ενέργειας και η
συμπεριφορά των μοντέλων σε συσκευές περιορισμένων πόρων δεν
αξιολογήθηκαν στην παρούσα εργασία.

\section{Μελλοντικές επεκτάσεις}
\label{sec:conclusions-future}

Τα αποτελέσματα της εργασίας αφήνουν ανοιχτές αρκετές κατευθύνσεις για
περαιτέρω διερεύνηση. Καθώς η αύξηση των δεδομένων απόσταξης βελτίωσε σημαντικά την επίδοση,
ιδιαίτερα για τα μικρότερα μοντέλα, μια άμεση επέκταση είναι η απόσταξη
με περισσότερα δεδομένα, είτε από σύνολα μεγαλύτερης κλίμακας, όπως το
\tl{Kinetics-400} \cite{kay2017kinetics}, είτε μέσω συλλογής βίντεο από
το διαδίκτυο (\tl{web scraping}). Παράλληλα, θα μπορούσαν να
εξεταστούν πιο σύνθετοι στόχοι απόσταξης, όπως οι αναπαραστάσεις
ενδιάμεσων επιπέδων του \tl{backbone} ή οι πίνακες αυτοπροσοχής
(\tl{self-attention}), κατά τα πρότυπα των \tl{ViTKD}
\cite{yang2024vitkd} και \tl{TinyMIM} \cite{ren2023tinymim}. Οι
στόχοι αυτοί μπορούν επίσης να συνδυαστούν με την απόσταξη των τελικών
προβλέψεων σε μια διευρυμένη υβριδική στρατηγική, με κατάλληλη
στάθμιση των επιμέρους όρων της συνάρτησης απώλειας.

Μια ακόμη κατεύθυνση αφορά τον συνδυασμό της απόσταξης γνώσης με τις
υπόλοιπες τεχνικές συμπίεσης που παρουσιάστηκαν στο
Κεφάλαιο~\ref{ch:sota}, ώστε να μειωθούν περαιτέρω οι απαιτήσεις
μνήμης και ο χρόνος εκτέλεσης. Ειδικότερα, τα αποσταγμένα μοντέλα-\tl{student}
θα μπορούσαν να κβαντοποιηθούν μετά την εκπαίδευσή τους με μεθόδους
σχεδιασμένες για αρχιτεκτονικές \tl{Vision Transformer}, όπως το
\tl{RepQ-ViT} \cite{li2023repqvit}. Για να αξιολογηθεί όχι μόνο η επίδραση της κβαντοποίησης
στην ακρίβεια αλλά και το πραγματικό όφελος σε ταχύτητα, μνήμη και
κατανάλωση ενέργειας, τα κβαντοποιημένα μοντέλα θα πρέπει να
εκτελεστούν σε \tl{hardware} που υποστηρίζει πράξεις χαμηλής
αριθμητικής ακρίβειας.

Τέλος, τα αποσταγμένα μοντέλα θα μπορούσαν να αξιολογηθούν σε
περισσότερα σύνολα δεδομένων και σε άλλες εργασίες κατανόησης βίντεο,
όπως η ανίχνευση συμβάντων και η παρακολούθηση αντικειμένων, ώστε να
εξεταστεί κατά πόσο οι αναπαραστάσεις που μαθαίνουν γενικεύουν πέρα
από την αναγνώριση ενεργειών. Ενδιαφέρον παρουσιάζει
επίσης η εφαρμογή της ίδιας μεθοδολογίας με νεότερα μοντέλα θεμελίωσης
ως \tl{teacher}, όπως το \tl{V-JEPA 2} \cite{assran2025vjepa2}.
